% 346_Ladungsquantisierung_Galois_En_ch.tex
% Johann Pascher, 4 September 2026

\providecommand{\BM}{\textbf{[B]}}
\providecommand{\KM}{\textbf{[K]}}
\providecommand{\SM}{\textbf{[S]}}

\phantomsection
\addcontentsline{toc}{chapter}{Charge Quantisation from GF$(27)^*$}

\section*{Status markers}
\begin{center}
\begin{tabular}{@{}lp{10.5cm}@{}}
\textbf{[K]} & \emph{Confirmed} --- algebraically complete. \\[4pt]
\textbf{[B]} & \emph{Established} --- numerically secured by script. \\[4pt]
\textbf{[S]} & \emph{Speculative} --- conjecture, not yet fully proved. \\
\end{tabular}
\end{center}

\section*{Abstract}

Electric charge quantisation and the colour singlet/triplet structure of the
Standard Model follow from two algebraic properties of GF$(27)^*$ orbits:
(1) quadratic residue mod~13 $\leftrightarrow$ electrically neutral;
(2) $N\bmod3$ encodes colour charge. Their combination classifies all Standard
Model fermions completely without free parameters. Theorem~E proposes Gell-Mann--Nishijima
$Q=I_3+Y/2$ as a consequence of the Witt-pair operator decomposition --- a precise
open question for a future document.

\section{Background}

\textbf{Doc.~336 \KM.} The injection $\iota(Q)=(3Q\bmod3,\lfloor Q\rfloor)$
shows $N\bmod3$ classifies $Q=N/3$ completely.

\textbf{Doc.~339 \BM.} Eight gluons $= 4\,\mathbb{Z}_{13}$-orbits $\times\,2\,\mathbb{Z}_2$-parities;
three elements per orbit correspond to three colours.

\textbf{Doc.~341 \BM.} $N\bmod3\in GF(3)$ classifies colour singlet ($N\equiv0$)
and colour triplet ($N\equiv1,2$).

\section{Theorem A --- Quadratic residues and electric neutrality}

\begin{theorem}[Electric neutrality $\leftrightarrow$ QR mod 13 \BM]
\begin{itemize}
\item $O_1=\{1,3,9\}$ and $O_3=\{4,10,12\}$: all elements are quadratic residues
  mod~13. These orbits carry $Q_\mathrm{el}=0$.
\item $O_2=\{2,5,6\}$ and $O_4=\{7,8,11\}$: all elements are quadratic
  non-residues mod~13. These orbits carry $Q_\mathrm{el}\neq0$.
\end{itemize}
\end{theorem}

\begin{proof}
$\mathrm{QR}_{13}=\{1,3,4,9,10,12\}$.
$O_1,O_3\subset\mathrm{QR}_{13}$; $O_2,O_4\subset\mathrm{NQR}_{13}$.
\end{proof}

The Majorana layer $\{O_1,O_3\}$ (Doc.~343) is exactly the QR layer.
The Dirac layer $\{O_2,O_4\}$ is exactly the NQR layer. The two algebraic
criteria --- Majorana/Dirac and QR/NQR --- coincide and mean the same thing:
QR $=$ electrically neutral \BM.

\section{Theorem B --- Complete fermion classification}

\begin{theorem}[Two Galois properties classify all fermions \BM]
\begin{center}
\begin{tabular}{llll}
\toprule
QR mod 13 & $N\bmod3$ & Fermion type & Example \\
\midrule
QR (neutral) & $\equiv0$ (singlet) & Neutrino & $\nu_e,\nu_\mu,\nu_\tau$ \\
QR (neutral) & $\equiv1,2$ (triplet) & --- & (forbidden) \\
NQR (charged) & $\equiv0$ (singlet) & Ch.\ lepton & $e,\mu,\tau$ \\
NQR (charged) & $\equiv1,2$ (triplet) & Quark & up, down, \ldots \\
\bottomrule
\end{tabular}
\end{center}
\end{theorem}

The row ``QR + triplet = forbidden'' is algebraically forced: no electrically
neutral coloured fermion exists in the Standard Model. This is a direct
consequence of the GF$(27)^*$ structure.

\section{Theorem C --- Charge quantisation from the Legendre symbol}

\[
  Q_\mathrm{el}=0 \;\Longleftrightarrow\; \left(\tfrac{k}{13}\right)=+1
  \;\text{ for all } k\in O_i.
\]
Algebraic derivation of $Q\in\{0,\pm\tfrac{1}{3},\pm\tfrac{2}{3},\pm1\}$
from Galois structure alone \BM.

\section{Theorem D --- Three generations from orbit elements}

Each primitive orbit has exactly three elements under Frobenius $k\mapsto3k\bmod26$.
The three elements correspond to three generations \SM:

\begin{center}
\begin{tabular}{llll}
\toprule
Orbit & Gen.~1 & Gen.~2 & Gen.~3 \\
\midrule
$O_1=\{1,3,9\}$ & $\nu_e$ & $\nu_\mu$ & $\nu_\tau$ \\
$O_3=\{4,10,12\}$ & $\bar\nu_e$ & $\bar\nu_\mu$ & $\bar\nu_\tau$ \\
$O_4=\{7,8,11\}$ & $e^-$ & $\mu^-$ & $\tau^-$ \\
$O_2=\{2,5,6\}$ & $u$ & $c$ & $t$ \\
\bottomrule
\end{tabular}
\end{center}

The mass hierarchy within each generation requires the Yukawa formula of Doc.~338.

\section{Theorem E --- Gell-Mann--Nishijima from Witt-pair operators (open) \SM}

\textbf{Motivation.} The three Witt-pair number operators $N_1,N_2,N_3$ from
Doc.~344 satisfy $N=N_1+N_2+N_3$ (the total number operator). In the Standard
Model, the Gell-Mann--Nishijima relation is:
\[
  Q = I_3 + \frac{Y}{2}, \qquad
  I_3 = \frac{1}{2}(N_1-N_2), \qquad
  Y = \frac{1}{3}(N_1+N_2-2N_3).
\]

\textbf{Check for leptons:} Neutrino: $N_1=N_2=N_3=0$ $\Rightarrow$ $I_3=0$,
$Y=0$, $Q=0$ $\checkmark$. Positron: $N_1=N_2=N_3=1$ $\Rightarrow$ $I_3=0$,
$Y=2$, $Q=1$ $\checkmark$.

\textbf{Why it is \SM.} For quarks, the decomposition of $N$ into $N_1,N_2,N_3$
is not unique without an additional assignment rule. The formula gives correct
$Q$ for leptonic states but requires an explicit $N_1,N_2,N_3$ assignment for
each quark state to be verified. This is a precise open question: can the
Witt-pair decomposition of GF$(27)^*$ orbits force $I_3$ and $Y$ algebraically?

\textbf{Note on Deepseek's analysis.} An external check (Deepseek, 4~September 2026)
proposed $Q(g^k)=(k\bmod3)/3$ element-wise. This is correct for the total
number operator $N(g^k)=k\bmod3$ from Doc.~341, combined with $Q=N/3$ from
Doc.~336. However, $N$ is not constant within an orbit (e.g.\ $O_1=\{1,3,9\}$
has $k\bmod3\in\{1,0,0\}$), so $Q$ is not a single value per orbit element
--- it is an orbit-averaged or ground-state property. The orbit-level property
(Theorem~A, QR$\leftrightarrow Q=0$) is the correct statement at orbit level.
Gell-Mann--Nishijima requires the individual $N_1,N_2,N_3$ decomposition and
remains \SM.

\section*{Verification script}

\texttt{pruef\_346\_ladungsquantisierung.py} --- Theorems A--C verified \BM.
Theorem~E: not yet scriptable (requires explicit $N_1,N_2,N_3$ assignment).

\section{Summary}

\begin{table}[H]\centering
\caption{Results Doc.~346}
\begin{tabular}{lll}\toprule
Theorem & Content & Status \\\midrule
A & QR mod 13 $\leftrightarrow$ electric neutrality & \BM \\
B & QR + $N\bmod3$ classifies all fermions & \BM \\
C & Charge quantisation from Legendre symbol & \BM \\
D & Three generations = three Frobenius powers & \SM \\
E & Gell-Mann--Nishijima from Witt-pair operators & \SM \\\bottomrule
\end{tabular}
\end{table}

Open: SU$(2)_L$ doublet structure explicitly from the Galois algebra.

\section*{Note on AI assistance}

Prepared with support from Claude (Anthropic). External algebraic check by Deepseek (4~September 2026).
Physical interpretations and conclusions by Johann Pascher (ORCID 0009-0000-6518-4064).

\begin{thebibliography}{9}
\bibitem{P336} J.~Pascher, \textit{Doc.~336: GF(9)-FFGFT bridge}, FFGFT corpus (Aug.\ 2026).
\bibitem{P339} J.~Pascher, \textit{Doc.~339: Frobenius separation}, FFGFT corpus (Aug.\ 2026).
\bibitem{P341} J.~Pascher, \textit{Doc.~341: GF(27) in GALG and FFGFT}, FFGFT corpus (Sept.\ 2026).
\bibitem{P343} J.~Pascher, \textit{Doc.~343: Zeta function of T$^4$/Z$_3$}, FFGFT corpus (Sept.\ 2026).
\bibitem{P344} J.~Pascher, \textit{Doc.~344: Furey fermions in GALG and FFGFT}, FFGFT corpus (Sept.\ 2026).
\bibitem{Furey2015} C.~Furey, Ph.D.\ Thesis, Waterloo 2015, arXiv:1611.09182.
\end{thebibliography}
